\chapter{Simple poles and higher collision operations}
\label{chap:sc050}

A singular product may become zero after taking cohomology. The
primitive that makes it zero can still have a singular product with a
third field. This gives a concrete way for a higher operation to survive
when the binary operation has disappeared.

The example below has only simple poles between its generating fields.
Its first surviving operation can nevertheless have any prescribed
arity greater than two. Reducing the fields to operators of energy
zero gives a two-by-two matrix algebra. The same polynomial controls
the higher operations of this matrix algebra and the insertions in its
Hochschild cochains. The comparison requires the fields, their
differential, the energy grading, and the signs of the higher maps.

\section{A differential algebra of free fields}
\label{sec:sc050-fields}

Fix an integer \(r\ge2\) and work over \(\mathbb C\). A parity is an
element of \(\mathbb Z/2\). Even quantities commute with scalars, and
interchanging two odd quantities in an exterior algebra changes the
sign. Consider the vector space
\[
 V_r=\mathbb C[x_0,x_1,\ldots]\otimes
       \Lambda(\gamma_0,\bar\gamma_0,\gamma_1,\bar\gamma_1,\ldots).
\]
The variables \(x_n\) are even. The variables \(\gamma_n,\bar\gamma_n\)
are odd, anticommute, and have square zero in this displayed
polynomial--exterior presentation. A state is a finite linear
combination of monomials. No completion of the state space is taken.

For an odd variable \(u\), the left derivative \(\partial_u\) is
specified by \(\partial_u u=1\) and
\[
 \partial_u(AB)=(\partial_u A)B+(-1)^{|A|}A\,\partial_u B
\]
on homogeneous monomials. Multiplication by \(u\) and \(\partial_u\)
have anticommutator one. Derivatives and multiplications associated
with distinct odd variables have anticommutator zero.

Put \(x=x_0\), \(\gamma=\gamma_0\), \(\bar\gamma=\bar\gamma_0\).
Define the following operator-valued formal series:
\begin{align}
 x(z)&=\sum_{n\ge0}x_nz^n, \notag\\
 \gamma(z)&=\sum_{n\ge0}\gamma_nz^n+
                    \sum_{n\ge0}\partial_{\bar\gamma_n}z^{-n-1},
       \label{eq:sc050-fields}\\
 \bar\gamma(z)&=\sum_{n\ge0}\bar\gamma_nz^n+
                    \sum_{n\ge0}\partial_{\gamma_n}z^{-n-1}.
       \notag
\end{align}
The coefficients \(x_n,\gamma_n,\bar\gamma_n\) in these formulas act
by multiplication. On each state, only finitely many derivatives act
nontrivially. A field therefore sends a state to a Laurent series
with a finite lower bound on its powers of \(z\).
The formal residue \(\operatorname{Res}_{z=0}\) extracts the
coefficient of \(z^{-1}\); it requires no analytic contour.

The variable \(z\) describes position. The variable \(x\) is a field
coordinate, not the position. This distinction matters when derivatives
or changes of coefficients are introduced.

The vacuum is \(1\). Translation is the even operator
\[
 T=\sum_{n\ge0}(n+1)
 \left(x_{n+1}\partial_{x_n}
       +\gamma_{n+1}\partial_{\gamma_n}
       +\bar\gamma_{n+1}\partial_{\bar\gamma_n}\right).
\]
Each sum is finite on a state. Direct computation gives
\([T,A(z)]=\partial_z A(z)\) for each generating field \(A\).
Also \(T1=0\) and \(T^n x/n!=x_n\), with the same identities for
the odd generators.

For homogeneous operators write
\([A,B]_{\mathrm{gr}}=AB-(-1)^{|A||B|}BA\).
The only nonzero generating commutator is
\[
 [\gamma(z),\bar\gamma(w)]_{\mathrm{gr}}
   =\delta(z,w)\,\mathrm{id},\qquad
 \delta(z,w)=\sum_{m\in\mathbb Z}z^{-m-1}w^m.
\]
The identity
\[
 \delta(z,w)=i_{z,w}(z-w)^{-1}-i_{w,z}(z-w)^{-1}
\]
uses expansion in powers of the second coordinate divided by the
first. Thus the singular products are
\begin{equation}\label{eq:sc050-poles}
 \gamma(z)\bar\gamma(w)\sim\frac1{z-w},\qquad
 \bar\gamma(z)\gamma(w)\sim\frac1{z-w}.
\end{equation}
The symbol \(\sim\) retains only negative powers of \(z-w\).
Every product involving \(x\) has no singular part.

\subsection{Composite states and normal products}

Normal ordering moves multiplication operators to the left of
derivatives, with a minus sign for each interchange of odd operators.
Associate to a state monomial the normally ordered product of the
corresponding divided derivatives of the fields in
\eqref{eq:sc050-fields}. Extend this assignment linearly and denote it
by \(Y(A,z)\). In particular \(Y(1,z)=\mathrm{id}\).
For states \(A,B\), define their products \(A_{(n)}B\) by
\[
 Y(A,z)B=\sum_{n\in\mathbb Z}(A_{(n)}B)z^{-n-1}.
\]
The normal product of two states is \(A_{(-1)}B\).
The vector-space presentation of \(V_r\) does not assert that every
normal product equals multiplication in its polynomial--exterior
coordinates.

Every product of composite fields can be calculated by moving the
derivatives in the first field past the multiplications in the
second. At a crossing, retain the reordered term or take the
anticommutator. Each odd factor can participate in at most one
contraction. A contraction between divided derivatives of orders
\(m,n\) of the two conjugate fields contributes
\begin{equation}\label{eq:sc050-contraction}
 \frac{\partial_z^m\partial_w^n(z-w)^{-1}}{m!\,n!}
 =\frac{(-1)^m(m+n)!}{m!\,n!}(z-w)^{-m-n-1}.
\end{equation}
The sign of a pairing is the sign of the permutation of odd factors
required to make its contractions. This gives a finite sum over
partial pairings for any fixed input monomials.

A vertex superalgebra is a parity-graded state space with a vacuum,
translation and a state--field assignment satisfying lower
truncation, the vacuum and translation identities, graded skew
symmetry, and the collision Jacobi identity. For the fields here,
the following proof states these identities and derives them
directly from the pairings.

\begin{proposition}\label{prop:sc050-vertex}
These fields give \(V_r\) a vertex superalgebra structure. They obey
the vacuum and translation identities
\[
 Y(A,z)1=e^{zT}A,\qquad
 Y(TA,z)=\partial_zY(A,z),\qquad
 [T,Y(A,z)]=\partial_zY(A,z),
\]
and the graded skew-symmetry identity
\[
 Y(A,z)B=(-1)^{|A||B|}e^{zT}Y(B,-z)A.
\]
Their two ordered products and their collision iterate satisfy the
vertex Jacobi identity.
\end{proposition}
\begin{proof}
On the vacuum all derivatives vanish. The remaining multiplication
series is the Taylor series \(e^{zT}A\). Differentiating a divided
derivative field gives the next one with its required integer
coefficient. This proves the two translation identities. Interchanging
the two lists of odd factors in the pairing expansion gives the
displayed skew sign. Formula \eqref{eq:sc050-contraction} transforms
with exactly the sign of the interchanged contracted factors.

For three states, expand their three fields into pairings. After
putting the third coordinate at zero, the result is a finite sum of
elements of
\[
 V_r[[u,v]][u^{-1},v^{-1},(u-v)^{-1}].
\]
The two radial expansions give the two ordered products, with the
graded sign when the first two states are exchanged. For the collision
expansion, first choose the pairs between those two states. Pair the
remaining factors with the third state and Taylor-expand their
uncontracted factors. This is a bijection with the original pairings.
It preserves the permutation of odd factors and hence every sign.
The collision expansion is therefore \(Y(Y(A,u-v)B,v)C\).

For a rational function of \(u\) with possible poles at \(0\) and
\(v\), partial fractions say that its outer residue is the sum of the
residues at these two poles. Taking the coefficient of \(v^{-1}\)
gives the difference between the two ordered expansions as the
collision residue. The same statement holds for the displayed formal
series, because only finitely many Taylor coefficients contribute at
fixed pole orders. Applying it after multiplication by arbitrary
Laurent monomials gives every coefficient of the vertex Jacobi
identity.
\end{proof}

A differential on a vertex superalgebra is an odd square-zero map
that commutes with translation and is a derivation of every product.
Let \(X(t)=\sum_{n\ge0}x_nt^n\). On the vector-space coordinates define
\begin{equation}\label{eq:sc050-Q}
 Qx_n=0,\qquad Q\gamma_n=x_n,\qquad
 Q\bar\gamma_n=\frac{[t^n]X(t)^r}{r+1},
\end{equation}
and extend as an odd derivation. Here \([t^n]\) means coefficient
extraction. Its square is zero on each generator, hence on every
monomial. Its commutators with the three fields are respectively
\[
 0,\qquad x(z),\qquad \frac{:x(z)^r:}{r+1}.
\]
These follow directly from the left derivative formulas. They
preserve the splitting into multiplication and derivative parts.
Thus \(Q\) is a derivation of every normally ordered composite and
every coefficient of its field. Formula \eqref{eq:sc050-Q} also
commutes with \(T\). It defines a differential on \(V_r\).

The closed odd field of interest is
\begin{equation}\label{eq:sc050-psi}
 \psi=\bar\gamma-\frac1{r+1}:x^{r-1}\gamma:.
\end{equation}
For states \(A,B\), the polynomial
\[
 [A_\lambda B]=\sum_{n\ge0}\frac{\lambda^n}{n!}A_{(n)}B
\]
records their singular products. It is called the \(\lambda\)-bracket.
Contraction of the fields in \eqref{eq:sc050-psi} gives
\begin{equation}\label{eq:sc050-primary}
 Q\psi=0,\qquad
 [\psi_\lambda\psi]=-\frac2{r+1}x^{r-1},\qquad
 [:x^j\gamma:_\lambda\psi]=x^j.
\end{equation}
The first bracket is exact, since
\(Q(:x^{r-2}\gamma:)=x^{r-1}\). Its primitive still has a singular
product with \(\psi\).

\subsection{Cohomology and its commutative product}

A state is closed if its differential is zero, and exact if it is
the differential of another state. Cohomology is the quotient of
closed states by exact states. A linear map commuting with the
differentials is a chain map; it is a quasi-isomorphism when its
induced map on cohomology is an isomorphism.

Introduce \(G(t)=\sum_{n\ge0}\gamma_nt^n\) and
\[
 \psi_n=\bar\gamma_n-\frac{[t^n](X(t)^{r-1}G(t))}{r+1}.
\]
This is an invertible triangular change of polynomial--exterior
coordinates: solve for each \(\bar\gamma_n\) by adding back the
displayed expression. In these coordinates,
\[
 Q\gamma_n=x_n,\qquad Qx_n=Q\psi_n=0.
\]
Give each \(x_n\) and \(\gamma_n\) pair degree one and each \(\psi_n\)
pair degree zero. On pair degree \(N>0\), put
\[
 h=\frac1N\sum_{n\ge0}\gamma_n\partial_{x_n};
\]
put \(h=0\) on pair degree zero. Each sum is finite on a state.
The anticommutator of \(Q\) with the numerator is multiplication by
\(N\). Thus
\[
 Qh+hQ=1-ip,
\]
where \(p\) sets the \(x_n,\gamma_n\) to zero and \(i\) includes the
remaining exterior vector space. In particular,
\begin{equation}\label{eq:sc050-cohom}
 H(V_r,Q)=\Lambda(\psi_0,\psi_1,\ldots),\qquad
 T\psi_n=(n+1)\psi_{n+1}.
\end{equation}
The singular products on this cohomology vanish. In a pairing of
surviving \(\psi\)-fields, every contraction introduces \(x^{r-1}\)
or one of its derivatives. Such factors are central and cannot be
removed by further contractions. Their images under \(p\) vanish.
The vertex product on \eqref{eq:sc050-cohom} is consequently the
commutative differential exterior product, with
\(Y(A,z)B=(e^{zT}A)B\). This conclusion concerns the cohomology
vertex algebra. It does not remove the singular operations of \(V_r\).

\section{A finite presentation of the singular operations}
\label{sec:sc050-presentation}

To retain the primitives of \eqref{eq:sc050-primary}, introduce a
central even symbol \(k\), even symbols \(c_1,\ldots,c_{r-1}\), and
odd symbols \(\psi,h_0,\ldots,h_{r-2}\). Set \(c_0=k\).
Let
\[
 L_r=\mathbb Ck\oplus\mathbb C[T]\psi
 \oplus\bigoplus_{j=0}^{r-2}\mathbb C[T]h_j
 \oplus\bigoplus_{j=1}^{r-1}\mathbb C[T]c_j,\qquad Tk=0.
\]
The notation \(\mathbb C[T]u\) means the free span of the derivatives
\(u,Tu,T^2u,\ldots\). Define
\begin{align}
 Qh_j&=c_{j+1},&Q\psi&=Qc_j=Qk=0,\notag\\
 [\psi_\lambda\psi]&=-\frac2{r+1}c_{r-1},&
 [h_j{}_\lambda\psi]&=[\psi_\lambda h_j]=c_j.
 \label{eq:sc050-L}
\end{align}
All other brackets of these generators are zero. Derivative inputs
are determined by
\[
 [TA_\lambda B]=-\lambda[A_\lambda B],\qquad
 [A_\lambda TB]=(T+\lambda)[A_\lambda B].
\]
These identities are called sesquilinearity. Together with graded
skew symmetry and
\[
 [A_\lambda[B_\mu C]]
 =[[A_\lambda B]_{\lambda+\mu}C]
   +(-1)^{|A||B|}[B_\mu[A_\lambda C]],
\]
they define a Lie conformal superalgebra. All the bracket values in
\eqref{eq:sc050-L} are central, so every term of this Jacobi identity
is zero on generators. Sesquilinearity extends the equality to their
derivatives. The same check proves that \(Q\) is a bracket derivation.

The assignments
\begin{equation}\label{eq:sc050-Lmap}
 k\longmapsto1,\qquad c_j\longmapsto x^j,\qquad
 h_j\longmapsto:x^j\gamma:,\qquad
 \psi\longmapsto\bar\gamma-\frac{:x^{r-1}\gamma:}{r+1}
\end{equation}
preserve the differential, translation and brackets. Their
\(\mathbb C[T]\)-linear independence can be checked in the free
coordinate space: the \(\bar\gamma_n\) coefficients separate the
derivatives of \(\psi\), and powers of \(x\) separate the remaining
families. Thus this presentation retains an actual subspace of fields.

\subsection{The vacuum enveloping algebra}

A mode of a field \(A\) is its coefficient \(A_{(m)}\) in
\(A(z)=\sum_m A_{(m)}z^{-m-1}\).
Form the Lie superalgebra of the modes of \(L_r\), with
\[
 [A_{(m)},B_{(n)}]_{\mathrm{gr}}=[A_\lambda B]_{(m+n)}
\]
for the constant primary brackets in \eqref{eq:sc050-L}, and
\((TA)_{(m)}=-mA_{(m-1)}\).
The formula on the right takes the indicated mode of the bracket
value. Specialize \(k_{(-1)}=1\) and \(k_{(m)}=0\) for \(m\ne-1\).
Let every nonnegative mode annihilate a vector \(1\), and let the
negative modes generate its induced module \(U_r\).

There is a concrete basis for this module. Order the negative modes.
Even central modes have arbitrary multiplicity, and each odd mode has
multiplicity at most one. Reorder two odd modes by their
anticommutator, which is central. The two reorderings of three odd
modes agree because their Jacobi expression is zero. For a repeated
odd letter \(u\), use \(u^2=[u,u]/2\); its overlaps agree because
\([u,u]\) is central. These reductions terminate by word length and
the number of inverted pairs.

Their linear independence has an explicit test. Over the polynomial
ring of negative central modes, write \(B_{ij}\) for the symmetric
central anticommutator of two negative odd modes. On the exterior
space with variables \(e_i\), represent these odd modes by
\[
 e_i\wedge+\frac12\sum_j B_{ij}\partial_{e_j}.
\]
The sum is finite on each exterior polynomial. Its anticommutators
are \(B_{ij}\). Applied to \(1\), an ordered word has its exterior
monomial as highest exterior-degree term. Descending in that degree
proves independence of the ordered words. Positive modes can be
moved to the right using the same central relations, so the vacuum
condition introduces no additional relation among negative words.

Define fields on the ordered states by taking successive normal
products of divided derivatives. Their generating contractions are
the central fields in \eqref{eq:sc050-L} divided by \(z-w\).
Their derivatives are still central. Moving modes through a normal
product therefore gives a finite pairing expansion, with the central
factors left at their contraction positions and then Taylor-expanded.
The two orderings and the collision expansion agree term by term
by the pairing argument in Proposition~\ref{prop:sc050-vertex}.
It proves locality and the vertex Jacobi identity for these fields.
Their value on the vacuum recovers the ordered state, so the
state--field assignment is well-defined and injective.
This constructs the vacuum enveloping vertex algebra \(U_r\).
The derivation \(Q\) of its mode relations descends to its vacuum
module and to every field product.

\begin{theorem}\label{thm:sc050-vertex-resolution}
The assignments \eqref{eq:sc050-Lmap} extend to a surjective
differential vertex-algebra map
\[
 q_r:U_r\longrightarrow V_r.
\]
It induces an isomorphism on cohomology, preserving translation and
every vertex product.
\end{theorem}
\begin{proof}
The mode assignments preserve all relations and the vacuum. The
construction by successive normal products then preserves every
composite field. They are surjective because \(h_0\) gives
\(\gamma\), \(c_1\) gives \(x\), and
\(\psi+:x^{r-1}\gamma:/(r+1)\) gives \(\bar\gamma\).

For homogeneous negative modes define symmetrization by
\[
 \operatorname{sym}(u_1\cdots u_n)
   =\frac1{n!}\sum_{\sigma\in S_n}
        \epsilon(\sigma;u)\,u_{\sigma(1)}\cdots u_{\sigma(n)}.
\]
Here \(S_n\) is the set of permutations of \(n\) letters, and
\(\epsilon(\sigma;u)\) is the product of the signs from
interchanging odd letters. Symmetrization is a vector-space isomorphism
from the symmetric algebra on even generators tensored with the
exterior algebra on odd generators. It is triangular with respect
to the ordered basis just proved. Since \(Q\) is a derivation of
the mode relations, it commutes with symmetrization.
In these coordinates the pairs
\[
 T^nh_j\longmapsto T^nc_{j+1}
\]
are contractible. The homotopy which replaces a \(T^nc_{j+1}\) by
\(T^nh_j\), sums over all such factors and divides by total pair
degree proves that only the exterior algebra on \(T^n\psi\)
survives. This homotopy and the projection commute with translation.

The target contraction in \eqref{eq:sc050-cohom} leaves exactly the
same exterior generators. After the map \(q_r\), a symmetrized
source word projects to its exterior word in those generators.
All contraction corrections contain an \(x_n\), and vanish under
the target projection. Thus \(H(q_r)\) is the identity on the
displayed exterior bases. It is an isomorphism, and preservation
of the vertex operations follows from the already constructed
strict vertex map.
\end{proof}

In particular the normal square of the source generator need not
vanish. If \(C=[\psi_\lambda\psi]\), graded skew symmetry gives
\begin{equation}\label{eq:sc050-normal-square}
 2:\psi\psi:=TC.
\end{equation}
The exterior presentation used for the cohomology computation is a
vector-space presentation by symmetrization. Equation
\eqref{eq:sc050-normal-square} records the additional normal-product
term before taking cohomology.

\section{The first higher conformal operation}
\label{sec:sc050-higher}

The cohomology of the smaller complex \(L_r\) is
\[
 H_r=\mathbb Ck\oplus\mathbb C[T]\theta.
\]
Here \(\theta\) is odd and \(k\) is even with \(Tk=0\).
Let \(i(\theta)=\psi\), let \(p\) kill the pairs \(h_j,c_{j+1}\),
and set \(h(c_{j+1})=h_j\), with \(h=0\) on the other generators.
Then \(Qh+hQ=1-ip\), \(ph=hi=h^2=0\).

Higher brackets retain the compatibility equations for these
primitives. Their \(n\)-input output lies in
\[
 H_r[\lambda_1,\ldots,\lambda_n]/
       (T+\lambda_1+\cdots+\lambda_n).
\]
Applying \(T\) to the \(j\)-th input multiplies the output by
\(-\lambda_j\). When an output is inserted into another operation,
its momentum is the sum of the momenta of its inputs. These rules
define translation-equivariant compositions. Graded skew symmetry
uses the sign of a permutation times its Koszul sign; on repeated
odd inputs the resulting operations are symmetric.

We use exponential coefficients. On a formal odd input \(a\),
the bracket and map series are respectively
\[
 \sum_n\frac{\ell_n(a,\ldots,a)}{n!},\qquad
 \sum_n\frac{F_n(a,\ldots,a)}{n!}.
\]
For ordinary higher brackets, the Jacobi equations say that
the induced odd operation on the symmetric coalgebra of the
inputs with parity reversed has square zero. Its degree-\(n\)
part consists of \(n\)-letter tensor words modulo the Koszul
permutation relations.
That coalgebra splits a word into every ordered pair of complementary
subsets, with its Koszul sign. An operation extends by applying it
to each subset. Consequently an identity of these extensions is
determined by its one-output component. For conformal operations,
carry the momentum labels through each subset and each insertion.
These finite identities with polynomial momentum coefficients
define the conformal higher Jacobi and morphism equations.
A higher model is minimal when its unary operation is zero.
A higher morphism is a quasi-isomorphism when its linear
component is a quasi-isomorphism of complexes.

\begin{theorem}\label{thm:sc050-first-operation}
There is a higher conformal model of \(L_r\) on \(H_r\) with the
single nonzero primary bracket
\begin{equation}\label{eq:sc050-ell}
 \ell_{r+1}(\theta,\ldots,\theta)=(-1)^r r!\,k.
\end{equation}
All higher brackets with \(k\) as an input are zero.
A higher conformal quasi-isomorphism to \(L_r\) has components
\begin{align}
 F_1(\theta)&=\psi,&F_1(k)&=k,\notag\\
 F_n(\theta,\ldots,\theta)
   &=\frac{(-1)^n n!}{r+1}h_{r-n}
       &&(2\le n\le r).
 \label{eq:sc050-F}
\end{align}
All other higher primary components vanish. Sesquilinearity determines
every derivative-input component.
\end{theorem}
\begin{proof}
Every composite of two source brackets has \(k\) as an input and
vanishes. Thus their higher Jacobi identities hold. Introduce an
even parameter \(t\) and the odd polynomial
\[
 F(t)=t\psi+\frac1{r+1}\sum_{n=2}^r(-t)^n h_{r-n}.
\]
From \eqref{eq:sc050-L},
\begin{align*}
 QF(t)&=\frac1{r+1}\sum_{n=2}^r(-t)^n c_{r-n+1},\\
 \frac12[F(t)_\lambda F(t)]
   &=-\frac{t^2}{r+1}c_{r-1}
     +\frac1{r+1}\sum_{n=2}^r(-1)^nt^{n+1}c_{r-n}.
\end{align*}
Terms of each degree from two through \(r\) cancel. The remaining
identity is
\begin{equation}\label{eq:sc050-LMC}
 QF(t)+\tfrac12[F(t)_\lambda F(t)]
      =\frac{(-1)^r}{r+1}t^{r+1}k.
\end{equation}
Its coefficient of \(t^n/n!\) is exactly the one-output morphism
identity on \(n\) copies of \(\theta\). A source bracket can only
be followed by \(F_1\), since higher components vanish on \(k\).
Inputs containing \(k\) give zero on both sides. Multiplying each
input by its polynomial in \(T\) amounts to multiplying each
identity by the specified polynomial in the \(-\lambda_j\).
This is compatible with the momentum sums at insertions, so it
proves all derivative-input identities.

An identity with several output blocks is a sum over partitions
of the input labels, with one of the proved identities in one
block. The other blocks carry the same component maps on both
sides. Their Koszul signs are the signs of the same subset
permutations. Thus these component equations prove the complete
higher morphism identities, with every momentum label retained.
Finally, \(F_1\) is the inclusion in the displayed contraction,
so it induces the cohomology isomorphism.
\end{proof}

For \(r=2\), write \(h=h_0,c=c_1\). The calculation reduces to
\[
 Qh=c,\qquad [\psi_\lambda\psi]=-\tfrac23c,\qquad
 [h_\lambda\psi]=k,\qquad F_2=\tfrac23h.
\]
The binary equation is \(QF_2+[\psi,\psi]=0\).
At three inputs its correction gives
\(\ell_3=3[\psi,F_2]=2k\).
For \(r=3\), the corresponding equations are
\[
 F_2=\tfrac12h_1,\qquad F_3=-\tfrac32h_0,\qquad
 QF_3+3[\psi,F_2]=0,\qquad \ell_4=4[\psi,F_3]=-6k.
\]
All the generating poles in both examples are simple. Pole order
and the first nonzero higher arity measure different operations.

An invertible higher change of coordinates between minimal models
has an invertible linear component. At the first nonzero arity,
its comparison equation contains only that linear component and
the bracket in \eqref{eq:sc050-ell}; every term involving a lower
bracket vanishes. Hence it cannot remove this first operation.
This statement concerns the stated conformal model and its
translation-equivariant maps.

\subsection{Composite fields retain the obstruction}

The contraction used in Theorem~\ref{thm:sc050-vertex-resolution}
also determines higher brackets on the full cohomology of \(U_r\).
The following construction specifies them, including composite
inputs. Use the symmetrized coordinates of that proof and their
maps \(i,p,h\). Introduce a formal odd element \(a\) of the
cohomology space, with a separate supercommuting coefficient for
each basis vector. A coefficient has the parity needed to make its
product with that vector odd. Complete only in the number of these
coefficients. Define recursively
\begin{equation}\label{eq:sc050-full-transfer}
 F(a)=i(a)-hR(a),\qquad
 R(a)=\tfrac12[F(a),F(a)],\qquad
 \ell(a)=pR(a).
\end{equation}
In these series a bracket sums over decompositions of the labelled
inputs. For conformal brackets, attach an independent momentum to
each input and add momenta at each insertion. A derivative of a
series means the sum of insertions in one labelled block, with
the Koszul sign. No momentum parameters are identified. Coefficient
extraction and polarization define every component of \(F\) and
\(\ell\). Each coefficient uses finitely many binary trees.
Translation-equivariance of the contraction makes these operations
sesquilinear.

Here is a direct verification of the higher identities. Put
\(B=QF+R\). Since \(Qh=1-ip-hQ\), and \(Qi=0\),
\[
 B=i\ell+h[QF,F]=i\ell+h[B,F].
\]
The last equality uses \([[F,F],F]=0\), the graded Jacobi identity
on the odd element \(F\). Let \(F_*\ell\) denote the derivative
of the formal series \(F\) in its odd vector-field direction
\(\ell\). Applying this odd derivative to \eqref{eq:sc050-full-transfer}
changes sign when it passes the odd map \(h\). It gives
\[
 F_*\ell=i\ell+h[F_*\ell,F].
\]
An equation \(X=i\ell+h[X,F]\) has a unique solution coefficient
by coefficient, because \(F\) has no constant term. Hence
\[
 QF+\tfrac12[F,F]=F_*\ell.
\]
The vector field \(Qb+[b,b]/2\) on formal odd elements of the target
has square zero: its terms cancel by \(Q^2=0\), the derivation
identity, and Jacobi. Applying the chain rule to the last displayed
identity therefore gives \(F_*(\ell^2)=0\). Terms with two first
derivatives cancel because the vector fields are odd. Since
\(pF=a\), its derivative satisfies \(pF_*=1\), and \(\ell^2=0\).
These coefficient equations prove all higher Jacobi and morphism
identities. The argument preserves the independent momentum labels
throughout, so it proves the conformal identities as well.

\begin{proposition}\label{prop:sc050-full-obstruction}
On the full cohomology \eqref{eq:sc050-cohom}, the transferred
brackets of arities \(2,\ldots,r\) vanish. At arity \(r+1\), the
value on \(r+1\) copies of \(\psi_0\) is \((-1)^r r!\,1\).
Thus this first higher operation cannot be removed by an
invertible change of minimal conformal coordinates.
\end{proposition}
\begin{proof}
Assign auxiliary weights \(j\) to \(c_j\), \(j+1\) to \(h_j\),
and zero to \(\psi,k\); derivatives have the same auxiliary
weight as their generators. For a monomial subtract the number
of \(h\)-factors from its auxiliary weight.

A contraction of two \(\psi\)-factors increases this difference by
\(r-1\). A contraction of \(h_j\) with \(\psi\) replaces \(h_j\)
by \(c_j\), and preserves it. Other primary contractions vanish.
Reordering and symmetrizing words use exactly these contractions.
The homotopy replaces one \(c_{j+1}\) by \(h_j\), so it decreases
the difference by one. Translation does not change it.

Every term of a bracket of two surviving composite states starts
with at least one contraction of two \(\psi\)-factors. In an
\(n\)-input tree for \(\ell\), there are \(n-2\) homotopies.
The difference in every term before projection is therefore at
least \(r-1-(n-2)\). For \(n\le r\) it is positive. Projection
keeps only exterior words in derivatives of \(\psi\), for which
the difference is zero. It kills every such term.

On primitive inputs, \eqref{eq:sc050-full-transfer} stays in
\(L_r\). Its homotopy agrees there with the linear contraction
used in Theorem~\ref{thm:sc050-first-operation}. Thus that theorem
gives the stated arity-\((r+1)\) value. The first-arity comparison
argument applies to the full minimal model just constructed.
\end{proof}

\section{Energy zero and the Zhu quotient}
\label{sec:sc050-zhu}

The singular fields and the matrix algebra have different inputs.
To relate them, retain only operators that preserve energy.
An energy grading of a vertex algebra is an even diagonal operator
\(H\) satisfying
\[
 [H,Y(A,z)]=z\partial_zY(A,z)+Y(HA,z).
\]
If \(HA=h_AA\), the number \(h_A\) is the weight of \(A\).
Translation raises it by one, and a nonzero product has weight
\(h_{A_{(n)}B}=h_A+h_B-n-1\).
An operator with this property is also called a Hamiltonian.
Its existence does not assert that it is a mode of any state.

Set \(a=1/(r+1)\). On \(V_r\) assign weights
\[
 h_x=h_\gamma=a,\qquad h_{\bar\gamma}=1-a.
\]
On divided derivatives add their derivative order, and on state
monomials add weights. The field formulas prove the Hamiltonian
identity. The differential preserves weight. The weights of the
generators in \(U_r\) are
\[
 h_\psi=\frac r{r+1},\qquad
 h_{h_j}=\frac{j+1}{r+1},\qquad
 h_{c_j}=\frac j{r+1},\qquad h_k=0.
\]
They preserve the relations, and \(q_r\) preserves this grading.

\subsection{An associative product from collision coefficients}

For a state \(A\) of weight \(h_A\), define deformed products by
\[
 A_{[n]}B=\sum_{j\ge0}\binom{h_A}{j}A_{(n+j)}B,\qquad
 \binom{h}{j}=\frac{h(h-1)\cdots(h-j+1)}{j!}.
\]
The empty product is one. The sum is finite on fixed states.
These are the coefficients of \((1+z)^{h_A}Y(A,z)\).
Extend them to nonhomogeneous states by linearity. Put
\[
 A*B=A_{[-1]}B,\qquad
 \{A,B\}_H=\sum_{j\ge0}\binom{h_A-1}{j}A_{(j)}B,
\]
and let \(O_H(V)\) be the linear span of all \(A_{[-2]}B\).
The braces in this paragraph denote this deformed binary bracket,
not Hochschild insertion operations.

\begin{lemma}\label{lem:sc050-zhu-quotient}
The subspace \(O_H(V)\) is a two-sided ideal for \(*\), and the
quotient \(Z_H(V)=V/O_H(V)\) is a unital associative algebra.
In the quotient,
\begin{equation}\label{eq:sc050-zhu-relations}
 [TA]=-h_A[A],\qquad
 [A]*[B]-(-1)^{|A||B|}[B]*[A]=[\{A,B\}_H].
\end{equation}
No integrality assumption on the weights is needed.
\end{lemma}
\begin{proof}
Write \(\nabla=T+H\). Differentiating the deformed fields gives
\[
 ((\nabla+n+1)A)_{[n]}B=-nA_{[n-1]}B,\qquad
 A_{[-2]}B=(\nabla A)*B.
\]
In particular every product with index at most \(-2\) lies in
\(O_H(V)\), by induction on its negative index. Taking \(B=1\)
gives \(\nabla V\subset O_H(V)\), and repeated translation gives
\[
 \frac{T^jA}{j!}\equiv\binom{-h_A}{j}A\pmod{\nabla V}.
\]

The coefficient form of the vertex Jacobi identity gives
\begin{align}
 (A*B)*C-A*(B*C)
  =\sum_{j\ge0}\bigl(
     A_{[-j-2]}(B_{[j]}C)
    +(-1)^{|A||B|}B_{[-j-2]}(A_{[j]}C)\bigr).
 \label{eq:sc050-zhu-associator}
\end{align}
To obtain it with these deformed coefficients, multiply the two
field factors by \((1+z)^{h_A}\) and \((1+w)^{h_B}\) in the
two product expansions. In their collision expansion use
\[
 (1+z)^{h_A}
  =\sum_{j\ge0}\binom{h_A}{j}(z-w)^j(1+w)^{h_A-j}.
\]
Taking the coefficients with two product indices \(-1\) gives
\eqref{eq:sc050-zhu-associator}; the two radial expansions give
its two summands and their graded sign.

The same coefficient identity, with one residue weighted by
\((1+z)^{h_A-1}\), gives the derivation formula
\[
 \{A,B_{[n]}C\}_H
   =\{\!A,B\!\}_H{}_{[n]}C
      +(-1)^{|A||B|}B_{[n]}\{A,C\}_H.
\]
It follows by expanding \((1+w+(z-w))^{h_A-1}\);
the coefficients combine by
\(\sum_i\binom{u}{i}\binom{v}{j-i}=\binom{u+v}{j}\).
This binomial identity follows by multiplying the two series
\((1+t)^u,(1+t)^v\).

Graded skew symmetry, followed by the displayed reduction of
\(T^j/j!\), gives
\[
 A*B-(-1)^{|A||B|}B*A-\{A,B\}_H\in\nabla V.
\]
Indeed the coefficient of \(A_{(j)}B\) after this reduction is
\(\binom{h_A}{j+1}-\binom{h_A-1}{j+1}
 =\binom{h_A-1}{j}\); the translation terms discarded at each
step lie in \(\nabla V\).

Apply \eqref{eq:sc050-zhu-associator} with \(A\) replaced by
\(\nabla A\). It shows that
\(((\nabla A)*B)*C\) differs from
\((\nabla A)*(B*C)\) by elements of \(O_H(V)\).
Hence \(O_H(V)\) is a right ideal. The bracket derivation formula
shows \(\{V,O_H(V)\}_H\subset O_H(V)\).
The skew formula then makes it a left ideal as well.
Equation \eqref{eq:sc050-zhu-associator} proves associativity in
the quotient. The vacuum formulas give \(1*A=A*1=A\).
All calculations use formal binomial coefficients, so rational
weights satisfy the same identities.
\end{proof}

The algebra \(Z_H(V)\) is the Hamiltonian-twisted Zhu algebra.
To interpret it through energy zero, consider modules with
twisted fields. A weight-\(h_A\) generating field in such a
module has the expansion
\[
 A(z)=\sum_{m\in\mathbb Z}\mathsf A_m z^{-m-h_A}.
\]
The commutator with energy is
\([H,\mathsf A_m]=-m\mathsf A_m\).
Rational powers are interpreted on a finite cover of the punctured
coordinate disk. The displayed monodromy under \(z\mapsto e^{2\pi i}z\) is
\(e^{-2\pi ih_A}\). Equivalently, for \(g=e^{2\pi iH}\), our
convention is \(Y(gA,z)=Y(A,e^{-2\pi i}z)\). This fixes the
sign of the twist. It respects products because their weight
differences are integers.
The index \(m=0\) preserves energy. Translation gives
\(\mathsf{(TA)}_0=-h_A\mathsf A_0\), as in
\eqref{eq:sc050-zhu-relations}.
The collision residue of two zero modes is the product \(*\).
Thus the quotient retains their multiplication, not just their
commutator.

\subsection{No additional zero-mode relations}

For the two vertex algebras at hand, it remains to prove that the
relations between their primary zero modes are defining relations.
Both algebras have the ordered mode basis already constructed.
Their primary singular brackets are central linear combinations
of generators or of the vacuum.

Here is a basis argument applicable to this situation. Filter states
by the number of primary generators in an ordered negative-mode
word. A contraction reduces that number. Replacing each \(T^nu\)
by \(\nabla^nu\), for a primary generator \(u\), is triangular in
derivative order, with diagonal coefficient one. Replacing an
ordered normal product by the right-associated \(*\)-product is
triangular in generator number, again with diagonal coefficient
one. Consequently the right-associated ordered \(*\)-words in
\(\nabla^nu\) form a basis.

Let \(W_+\) be the span of those words that contain a factor with
\(n>0\). We claim \(W_+=O_H(V)\). The operator \(\nabla\) is an
exact derivation of all the deformed products. In fact
\[
 [\nabla,(1+z)^{h_A}Y(A,z)]
   =(1+z)\partial_z\bigl((1+z)^{h_A}Y(A,z)\bigr)
   =Y_H(\nabla A,z),
\]
where \(Y_H(A,z)=(1+z)^{h_A}Y(A,z)\).
Also
\[
 \{\nabla A,B\}_H=0,\qquad
 \{A,\nabla B\}_H=\nabla\{A,B\}_H.
\]
These follow by integration by parts in the defining formal
residue for \(\{-,-\}_H\).

Induction on total generator number shows that \(W_+\) is a
two-sided ideal. In re-associating a product of basis words, every
error in \eqref{eq:sc050-zhu-associator} has smaller generator
number, because it contains a singular contraction. Moreover,
\[
 A_{[-j-2]}B
   =\frac1{(j+1)!}
       \bigl(\nabla(\nabla-1)\cdots(\nabla-j)A\bigr)*B
\]
has a factor of \(\nabla\). The induction hypothesis places each
such error in \(W_+\). For the leading re-associated words, swap
adjacent generators into order. A swap involving \(\nabla^nu\),
\(n>0\), has zero deformed bracket when it is the first input;
when it is second, that bracket still contains a factor of
\(\nabla\). The remaining skew correction lies in \(\nabla V\)
and has smaller generator number. The derivation rule expands it
into words with positive derivative order, to which induction
applies. Odd-square reductions are the same skew calculation
divided by two. These steps prove the ideal assertion, starting
from the constants, whose \(\nabla\) is zero.

The derivation rule gives \(\nabla V\subset W_+\), so
\(O_H(V)=(\nabla V)*V\subset W_+\). Conversely every positive
factor \(\nabla^nu\) belongs to \(O_H(V)\), and that subspace is
an ideal by Lemma~\ref{lem:sc050-zhu-quotient}.
Thus \(W_+=O_H(V)\).

It follows that the ordered \(*\)-words in the primary generators,
without their positive derivatives, form a basis of \(Z_H(V)\).
Their graded commutator is their primary bracket, because all
primary brackets are independent of \(\lambda\).
There are no further relations.

For \(V_r\) this proves
\begin{equation}\label{eq:sc050-clifford}
 Z_H(V_r)=
 \mathbb C[x]\langle\gamma,\bar\gamma\rangle/
   (\gamma^2,\bar\gamma^2,\gamma\bar\gamma+\bar\gamma\gamma-1),
\end{equation}
with central \(x\). The two scalar corrections to the mixed
products are \(h_\gamma\) and \(h_{\bar\gamma}\); their sum is one.
For \(U_r\), let \(\mathfrak l_r\) be the finite Lie superalgebra
on the primary generators in \eqref{eq:sc050-L}, omitting \(T\).
Then
\begin{equation}\label{eq:sc050-B}
 Z_H(U_r)=B_r:=U(\mathfrak l_r)/(k-1).
\end{equation}
Here the enveloping algebra is the tensor algebra divided by
\(AB-(-1)^{|A||B|}BA-[A,B]\).
For the odd generator \(\psi\), equation
\eqref{eq:sc050-normal-square} gives the same odd-square relation:
if \(C=[\psi_\lambda\psi]\), then
\[
 [\psi]*[\psi]=(-h_C/2+h_\psi)[C]=[C]/2,
 \qquad h_C=2h_\psi-1.
\]

\section{The matrix differential and the comparison square}
\label{sec:sc050-matrix}

Let \(M=\mathbb C[x]^{1|1}\), the free module with one even and
one odd basis vector. Its even endomorphisms are diagonal matrices;
its odd endomorphisms are off-diagonal matrices. Set
\[
 W(x)=\frac{x^{r+1}}{r+1},\qquad
 D=\begin{pmatrix}0&x\\x^r/(r+1)&0\end{pmatrix}.
\]
A matrix factorization of a scalar \(W\) is such an odd operator
with square \(W\,\mathrm{id}\). Here \(D^2=W(x)I\).
Its endomorphism algebra \(A_r=\operatorname{End}_{\mathbb C[x]}M\)
has differential
\[
 \delta A=DA-(-1)^{|A|}AD.
\]
Expansion gives \(\delta^2A=D^2A-AD^2=0\), and gives the graded
Leibniz identity for multiplication.

\begin{theorem}\label{thm:sc050-zhu-matrix}
The assignments
\[
 x\longmapsto xI,\qquad
 \gamma\longmapsto E_{21},\qquad
 \bar\gamma\longmapsto E_{12}
\]
give an isomorphism of differential algebras \(Z_H(V_r)\cong A_r\).
There is a commuting diagram
\begin{equation}\label{eq:sc050-square}
 \begin{array}{ccc}
 U_r&\xrightarrow{\ q_r\ }&V_r\\
 \big\downarrow\pi_U&&\big\downarrow\pi_V\\
 B_r&\xrightarrow{\ \rho_r\ }&A_r ,
 \end{array}
\end{equation}
where the upper arrow is a differential vertex quasi-isomorphism
and the lower arrow is a differential associative quasi-isomorphism.
The vertical arrows are the quotient maps of complexes, with
the matrix identification on the right. Their targets retain
the deformed product \(*\); they are not asserted equivalences.
\end{theorem}
\begin{proof}
The four ordered words \(1,\gamma,\bar\gamma,\gamma\bar\gamma\)
in \eqref{eq:sc050-clifford} are a free \(\mathbb C[x]\)-basis.
Their matrices are \(I,E_{21},E_{12},E_{22}\), which form a basis
of the matrix algebra. The relations hold by multiplication.
The map is therefore an algebra isomorphism. Since \(Q\) commutes
with \(H,T\), it preserves the deformed products and their ideal.
The differential descends, and direct multiplication gives
\[
 \delta E_{21}=xI,\qquad
 \delta E_{12}=\frac{x^r}{r+1}I.
\]
These are exactly its two generating differentials.

The map on \(B_r\) sends
\[
 c_j\longmapsto x^jI,\qquad
 h_j\longmapsto x^jE_{21},\qquad
 \psi\longmapsto
 \vartheta=E_{12}-\frac{x^{r-1}}{r+1}E_{21}.
\]
The differential and bracket relations follow by direct
supercommutation. This is the map induced from \(q_r\), proving
commutativity of the diagram. It is onto, since \(h_0,c_1\) give
\(E_{21},xI\), and
\(\psi+c_{r-1}h_0/(r+1)\) gives \(E_{12}\).

The ordered basis of \(B_r\), or its symmetrization, identifies
its complex with
\[
 \mathbb C[c_1,\ldots,c_{r-1}]
 \otimes\Lambda(h_0,\ldots,h_{r-2},\psi),
\]
with \(Qh_j=c_{j+1}\). The pair-degree contraction leaves
cohomology with basis \(1,\psi\).
For the target, a diagonal matrix satisfies
\[
 \delta\begin{pmatrix}a&0\\0&d\end{pmatrix}
      =(d-a)x\vartheta,
\]
and an odd matrix satisfies
\[
 \delta(bE_{12}+cE_{21})
      =x\left(c+\frac{x^{r-1}b}{r+1}\right)I.
\]
Since \(x\) is not a zero divisor, the even cycles are
\(\mathbb C[x]I\) and the odd cycles are
\(\mathbb C[x]\vartheta\).
Their boundaries are respectively \(x\mathbb C[x]I\) and
\(x\mathbb C[x]\vartheta\).
Thus \(H(A_r)\) has basis \(I,\vartheta\), and \(\rho_r\)
induces the identity on the two displayed cohomology bases.
The upper quasi-isomorphism was proved independently in
Theorem~\ref{thm:sc050-vertex-resolution}.
\end{proof}

The same formulas restricted to \(\mathfrak l_r\) give a
differential Lie quasi-isomorphism to \(A_r\) with its
supercommutator. Its source cohomology consists of \(k,\psi\);
its target cohomology was computed in the proof.
The finite Lie map, the associative map, and the vertex map
therefore have separate proofs at their stated operation levels.

Completing \(\mathbb C[x]\) at \(x=0\) gives the corresponding
formal matrix algebra over \(\mathbb C[[x]]\).
The kernel and image calculations above remain valid there,
because multiplication by \(x\) is still injective.
This is a specified coefficient change after the polynomial
comparison. No assertion about commuting an arbitrary Zhu
reduction with completion or cohomology is required.

\subsection{A twisted Fock module with a curved charge}

The matrices above arise on the energy-zero states of an explicit
field module. Introduce even variables \(X,X_{-1},X_{-2},\ldots\)
and odd exterior variables \(G_{-n},\bar G_{-n}\), \(n\ge1\).
Let \(S_0=\mathbb C^{1|1}\) and put
\[
 \mathcal M=
 \mathbb C[X,X_{-1},X_{-2},\ldots]\otimes
 \Lambda(G_{-1},\bar G_{-1},G_{-2},\bar G_{-2},\ldots)
 \otimes S_0.
\]
Give \(X\) and \(S_0\) energy zero, and each variable with
subscript \(-n\) energy \(n\). Energies add on monomials.
The grading operator \(\mathcal H\) multiplies a monomial by
its energy. All states again have finite support.

For \(m<0\), let \(X_m,\Gamma_m,\bar\Gamma_m\) act by
multiplication by \(X_m,G_m,\bar G_m\).
Set \(X_0=X\), \(X_m=0\) for \(m>0\), and
\[
 \Gamma_m=\partial_{\bar G_{-m}},\qquad
 \bar\Gamma_m=\partial_{G_{-m}}\quad(m>0).
\]
On \(S_0\) let \(\Gamma_0=E_{21}\) and
\(\bar\Gamma_0=E_{12}\). These odd zero modes act on the tensor
product with the sign
\[
 \Gamma_0(A\otimes v)=(-1)^{|A|}A\otimes E_{21}v,
\]
and similarly for \(\bar\Gamma_0\). Therefore all mode relations are
\[
 \{\Gamma_m,\bar\Gamma_n\}=\delta_{m+n,0}I,\qquad
 \{\Gamma_m,\Gamma_n\}=\{\bar\Gamma_m,\bar\Gamma_n\}=0,
\]
with all \(X_m\) central. Each mode of index \(m\) has
commutator \(-m\) times itself with \(\mathcal H\).

With \(a=1/(r+1)\), the twisted generating fields are
\begin{align*}
 x^{\mathcal M}(z)&=\sum_{m\le0}X_mz^{-m-a},\\
 \gamma^{\mathcal M}(z)&=\sum_{m\in\mathbb Z}\Gamma_mz^{-m-a},\\
 \bar\gamma^{\mathcal M}(z)&=
       \sum_{m\in\mathbb Z}\bar\Gamma_mz^{-m-(1-a)}.
\end{align*}
Their annihilation parts act finitely on each state.
Their contractions have precisely the required local pole.
For example, positive \(\Gamma\)-modes contracted with negative
\(\bar\Gamma\)-modes give
\begin{equation}\label{eq:sc050-twisted-contraction}
 \sum_{m>0}z^{-m-a}w^{m-(1-a)}
       =\frac{z^{-a}w^a}{z-w}.
\end{equation}
Near \(z=w\) its singular part is \((z-w)^{-1}\).
Define the fields of composite states by expanding at \(z=w\)
and subtracting the singular part prescribed by
\eqref{eq:sc050-poles}, then taking the required regular
coefficient. This is point-splitting normal ordering. For a
derivative state, differentiate the corresponding field.
The binomial expansion of \(z^{-a}=w^{-a}(1+(z-w)/w)^{-a}\)
specifies every correction.

These rules give a twisted module for \(V_r\). Indeed, each
multiple product is a finite pairing sum, with the factors in
\eqref{eq:sc050-twisted-contraction} and their conjugates and
derivatives. Expanding these same factors in the two radial
orders gives the two field products. Expanding them first at a
collision gives the stated point-splitting definition.
The bijection of pairings preserves the graded sign.
Formal residues therefore prove the twisted product and
iterate identities exactly as in
Proposition~\ref{prop:sc050-vertex}.
The only added factors are the displayed rational powers, whose
binomial expansions have just fixed the monodromy.

The normal-ordering correction can already be seen on
\(\mathcal M_0=\mathbb C[X]\otimes S_0\), the energy-zero subspace.
The regular part of \eqref{eq:sc050-twisted-contraction} at
\(z=w\) is \(-a/w\). The zero modes contribute
\(\Gamma_0\bar\Gamma_0/w\). Hence
\[
 (: \gamma\bar\gamma :)_0
       =\Gamma_0\bar\Gamma_0-aI
       \quad\hbox{on }\mathcal M_0.
\]
Adding the scalar correction \(aI\) in the Zhu product recovers
\(\Gamma_0\bar\Gamma_0\). This is the product in
\eqref{eq:sc050-clifford} acting on the stated module.

Define an odd energy-preserving operator on \(\mathcal M\) by
\begin{equation}\label{eq:sc050-module-charge}
 \mathcal Q=
   \sum_{m\le0}X_m\bar\Gamma_{-m}
   +\frac1{r+1}
      \sum_{m_1,\ldots,m_r\le0}
          X_{m_1}\cdots X_{m_r}\Gamma_{-(m_1+\cdots+m_r)}.
\end{equation}
On a fixed state, positive fermion modes beyond its largest
occupied index vanish. For each remaining index there are
finitely many tuples with the prescribed sum. Thus both sums
define algebraic operators on all finite states.
The mode anticommutators give
\[
 [\mathcal Q,\Gamma_m]_{\mathrm{gr}}=X_m,\qquad
 [\mathcal Q,\bar\Gamma_m]_{\mathrm{gr}}
   =\frac1{r+1}\sum_{m_1+\cdots+m_r=m}
                         X_{m_1}\cdots X_{m_r}.
\]
The sum uses \(m_i\le0\) and is zero for \(m>0\).
Also \([\mathcal Q,X_m]=0\). These are exactly the generating
field equations for \(Q\), so normal products and derivatives
give \([\mathcal Q,Y^{\mathcal M}(A,z)]_{\mathrm{gr}}
=Y^{\mathcal M}(QA,z)\) for every state \(A\).

The square of \(\mathcal Q\) is especially restrictive.
Only a cross anticommutator between the two sums in
\eqref{eq:sc050-module-charge} can survive. Both its fermion
indices are nonnegative, so their sum can be zero only when
both vanish. All the associated \(m,m_i\) are then zero.
Consequently
\begin{equation}\label{eq:sc050-module-curvature}
 \mathcal Q^2=\frac{X^{r+1}}{r+1}I=W(X)I.
\end{equation}
On \(\mathcal M_0\) the charge is
\[
 \mathcal Q|_{\mathcal M_0}
       =X E_{12}+\frac{X^r}{r+1}E_{21}=D(X).
\]
Thus the energy-zero module with its charge is the displayed
matrix factorization. It is a curved module: its charge has
the specified nonzero scalar square. Its endomorphism
supercommutator is an ordinary square-zero differential,
because that scalar commutes with every endomorphism.
This distinction is forced by the field representation itself.

\section{The higher comparison and its coefficient signs}
\label{sec:sc050-associative}

The higher conformal map also respects the energy-zero sector.
For a primary field of weight \(h\), its energy-zero vertex index
is \(h-1\). A constant conformal operation adds its input indices.
For the component \(F_n\) of \eqref{eq:sc050-F},
\[
 n h_\psi-(n-1)=\frac{r-n+1}{r+1}=h_{h_{r-n}}.
\]
Its output is therefore again of energy zero. At \(n=r+1\),
the sum of the input indices is \(-1\), the only nonzero mode
index of the constant field \(k\). Applying the matrix map gives
\begin{equation}\label{eq:sc050-matrix-F}
 F_1=\vartheta,\qquad
 F_n=\frac{(-1)^nn!}{r+1}x^{r-n}E_{21}\quad(2\le n\le r).
\end{equation}
Derivative inputs have the falling factorials from
\((TA)_{(m)}=-mA_{(m-1)}\). The binomial rule for differentiating
a product proves that this mode conversion commutes with
insertions, including their summed indices.

To compare with ordered multiplication, define the diagonal
matrix over \(\mathbb C[x,y]\)
\[
 \Delta(x,y)=
 \begin{pmatrix}
 0&x-y\\
 \dfrac1{r+1}\displaystyle\sum_{j=0}^r x^{r-j}y^j&0
 \end{pmatrix}.
\]
Multiplying the finite geometric sum by \(x-y\) gives
\(\Delta(x,y)^2=(W(x)-W(y))I\), and \(\Delta(x,0)=D\).
Thus the ordinary matrix series
\[
 L(t)=\Delta(x,-t)-D
     =t\vartheta+\frac1{r+1}
                         \sum_{n=2}^r(-t)^n x^{r-n}E_{21}
\]
satisfies
\begin{equation}\label{eq:sc050-Lmatrix}
 \delta L+L^2=-W(-t)I.
\end{equation}
Its exponential coefficients are exactly
\eqref{eq:sc050-matrix-F}. This is the matrix form of the
conformal identity \eqref{eq:sc050-LMC}.

\subsection{Ordered words and the boundary algebra}

The associative higher model uses ordered tensor words. For a
parity-graded vector space \(C\), let \(sC\) have the opposite
parity, and set
\[
 T^c(sC)=\bigoplus_{n\ge0}(sC)^{\otimes n}.
\]
Its coproduct cuts a word at every position. An odd operation
\(b_n:(sC)^{\otimes n}\to sC\) extends to words by replacing
each consecutive block of length \(n\); the sign is
\((-1)^{\text{parity of the preceding letters}}\).
An \(A_\infty\)-algebra is a family of such operations whose
extended sum \(b\) satisfies \(b^2=0\), with no zero-input
component in the algebras used here.

A morphism has even components \(\Phi_n\). Its extension replaces
each ordered partition into consecutive blocks by the tensor of
the corresponding component values. It is an \(A_\infty\)-map
when \(b'\Phi=\Phi b\). It is a quasi-isomorphism when \(\Phi_1\)
induces an isomorphism on cohomology. For a differential algebra
\((A,\delta)\), the convention is
\begin{equation}\label{eq:sc050-suspension}
 b_1(sA)=s(\delta A),\qquad
 b_2(sA,sB)=(-1)^{|A|}s(AB).
\end{equation}
These signs determine the comparison below.

Let \(H_W=\mathbb C1\oplus\mathbb C\theta\), with \(\theta\) odd.
Write \(e=s1\), which is odd, and \(a=s\theta\), which is even.
Set
\begin{equation}\label{eq:sc050-boundary}
 b_n(a^{\otimes n})=W_n e,\qquad
 W(t)=\sum_{n\ge3}W_nt^n,
\end{equation}
and add the unit operations
\[
 b_2(e,e)=e,\qquad b_2(e,a)=a,\qquad b_2(a,e)=-a.
\]
All other unit-containing operations vanish. These formulas
define a strictly unital algebra: after removing suspensions,
the even element \(1\) multiplies as an identity, has zero
differential, and kills every higher operation in which it
appears.

Here is a verification of \(b^2=0\) on all words. Its square is
a coderivation, since the two cross terms in the tensor
coproduct cancel for the odd map \(b\). It is enough to check
the one-output component. On a pure word \(a^N\), the only
two-step contributions replace \(N-1\) letters by \(W_{N-1}e\)
at one endpoint or the other. The final unit multiplication
gives \(W_{N-1}a-W_{N-1}a=0\).
For a single interior \(e\), removing it on its left or right
has opposite signs and leaves the same pure word.
For an endpoint \(e\), removing the unit and then applying the
pure operation cancels applying that operation first: its
prefix sign is negative at a left endpoint, and
\(b_2(a,e)=-a\) gives the negative sign at a right endpoint.
With at least two units, only three-letter words can reach one
output in two steps. The words \(eee,eea,eae,aee\) give pairs
\((e,-e),(a,-a),(-a,a),(a,-a)\).
The remaining words give zero. This exhausts the possibilities.

\begin{theorem}\label{thm:sc050-boundary-matrix}
For \(W(t)=t^{r+1}/(r+1)\), a strictly unital
\(A_\infty\)-quasi-isomorphism \(H_W\to A_r\) is given by
\[
 \Phi_1(e)=sI,\qquad
 \Phi_n(a^{\otimes n})
       =-\frac{s\,\partial_y^n\Delta(x,0)}{n!}\quad(n\ge1),
\]
with all other higher components zero.
\end{theorem}
\begin{proof}
Let \(K(t)=-(\Delta(x,t)-D)=\sum_{n\ge1}K_nt^n\).
Subtracting \(D^2\) from \((D-K)^2\) gives
\begin{equation}\label{eq:sc050-Kmatrix}
 \delta K-K^2=W(t)I.
\end{equation}
By \eqref{eq:sc050-suspension}, the one-output component of
\(b^{A_r}\Phi\) on \(a^n\) is
\[
 s\left(\delta K_n-\sum_{j=1}^{n-1}K_jK_{n-j}\right).
\]
The corresponding component of \(\Phi b\) is \(W_nsI\).
Any shorter scalar-output block would leave a unit in a word
of length at least two, on which the higher component is zero.
Thus \eqref{eq:sc050-Kmatrix} proves the pure-input equations.
On a word with one interior unit, the two adjacent removals
cancel. At an endpoint, multiplication by \(I\) on the target
gives respectively \(sK_n\) and \(-sK_n\), matching the two
source unit signs. The word \(ee\) gives \(sI\); words with more
units or letters have zero one-output difference.

The difference \(b^{A_r}\Phi-\Phi b\) is a coderivation along
\(\Phi\). Induction on word length with the cut coproduct
therefore extends the one-output equations to the full identity.
Finally \(K_1=\vartheta\), so the linear component induces the
basis \(I,\vartheta\) of the matrix cohomology computed in
Theorem~\ref{thm:sc050-zhu-matrix}. This proves the
quasi-isomorphism.
\end{proof}

The two series are related by \(L(t)=-K(-t)\). Hence the
exponential Lie component and the ordinary associative
coefficient satisfy
\begin{equation}\label{eq:sc050-sign-comparison}
 F_n^{\mathrm{Lie}}=(-1)^{n+1}n!K_n.
\end{equation}
For the cubic, \(K_2=-E_{21}/3\) whereas \(F_2^{\mathrm{Lie}}=2E_{21}/3\).
For the quartic, \(K_2=-xE_{21}/4\), \(K_3=-E_{21}/4\), while
\(F_2^{\mathrm{Lie}}=xE_{21}/2\), \(F_3^{\mathrm{Lie}}=-3E_{21}/2\).
The factorial and the change of parameter both enter the map.
Using \(L\) instead of \(-L\) as the associative series for
the reflected potential \(W(-t)\) leaves an arity-two error
\(2x^{r-1}I/(r+1)\). The two matrix identities fix the sign.

There is no strictly unital \(A_\infty\)-map from \(H_W\) to
\(\mathbb C\) which kills \(\theta\) and preserves these operations.
Every higher pure-input component would vanish by parity. Its
equation on \(r+1\) copies of \(\theta\) would then send the
nonzero scalar output to zero. Deleting that scalar term before
forming ordered words therefore changes the algebra. The
scalar polynomial must remain in any dual description.

\section{All normalized Hochschild insertions}
\label{sec:sc050-hochschild}

The ordered boundary model is small, but its operations act on
words of every length. A cochain records how such a word changes.
This is the source of the larger algebra of compatible boundary
operations.

Allow any series \(P(t)\in t^3\mathbb C[[t]]\), and define \(H_P\)
by \eqref{eq:sc050-boundary} with \(P\) in place of \(W\), retaining
the same units. The proof of \(b^2=0\) applies coefficient by
coefficient. A normalized Hochschild cochain is a multilinear
map on each input length, zero whenever a nonempty input word
contains the unit. Cochains with zero inputs are included.
Take a product, rather than a direct sum, over the input lengths.

In suspended notation the normalized cochains have basis
\[
 u_n(a^{\otimes n})=e,\qquad v_n(a^{\otimes n})=a,\qquad n\ge0.
\]
Their parities are one and zero respectively. For formal series
\(f(t)=\sum f_nt^n\), \(g(t)=\sum g_nt^n\), write
\[
 u(f)=\sum_{n\ge0}f_nu_n,\qquad
 v(g)=\sum_{n\ge0}g_nv_n.
\]
After removing the cochain output suspension, \(u(f)\) has even
parity and \(v(g)\) has odd parity.

The insertion brace \(F\{G_1,\ldots,G_k\}\) inserts the ordered
outputs of the \(G_i\) into distinct slots of \(F\), summing all
placements with the Koszul signs. Slots not used for an insertion
receive the remaining inputs in order. A zero-input insertion
occupies a slot while consuming no external input.
Define the Hochschild differential by
\[
 D F=[b,F]=b\{F\}-(-1)^{|F|}F\{b\}.
\]
Here parities are suspended parities, so \(|b|=1\).

\begin{theorem}\label{thm:sc050-braces}
For \(k\ge1\), all normalized insertion braces are
\begin{align}
 u(f)\{v(g_1),\ldots,v(g_k)\}
     &=u\left(\frac{f^{(k)}}{k!}\prod_i g_i\right),\notag\\
 v(f)\{v(g_1),\ldots,v(g_k)\}
     &=v\left(\frac{f^{(k)}}{k!}\prod_i g_i\right).
 \label{eq:sc050-braces}
\end{align}
An insertion of type \(u\) gives zero. The differential is
\begin{equation}\label{eq:sc050-HHd}
 Du(f)=0,\qquad Dv(g)=u(P'g).
\end{equation}
These formulas include zero-input cochains and the product
completion in input length.
\end{theorem}
\begin{proof}
Take \(f=t^m\) and \(g_i=t^{n_i}\).
For a nonzero insertion, each inserted output must be \(a\);
an \(e\) input kills a normalized outer cochain. Hence every
inserted cochain has type \(v\).
There are \(\binom{m}{k}\) ordered choices of \(k\) slots.
Each gives output length \(m-k+\sum_i n_i\), with positive
sign because the inserted maps and all external \(a\)'s are
even. This is the coefficient in
\(f^{(k)}g_1\cdots g_k/k!\). It proves both brace formulas,
including \(n_i=0\).

Split \(b\) into its scalar-output part \(u(P)\) and its binary
unit operations. The first part brackets with \(v(g)\) to give
\(u(P'g)\), by the brace formula, and brackets with \(u(f)\)
to give zero. On pure input words the binary unit contributions
cancel at the two endpoints. In the reverse insertion, the
binary product of two \(a\)'s is zero.
On words containing a unit, the two adjacent unit removals
cancel against the corresponding endpoint insertion. For a
zero-input cochain, the same two endpoint terms are the left
and right multiplication of its output, and have opposite
signs. Thus the normalized subspace is preserved and
\eqref{eq:sc050-HHd} holds in every input length.

At output length \(N\), a brace with \(k\) inserted cochains
uses outer length at most \(N+k\). There are finitely many
nonnegative input lengths summing to the required value.
Consequently each coefficient is a finite sum; the displayed
operations are defined on the product completion.
\end{proof}

The higher cup operations in suspended form are
\(M_1=D\) and \(M_k(F_1,\ldots,F_k)=b\{F_1,\ldots,F_k\}\)
for \(k\ge2\). The binary unit contributions are
\begin{align*}
 M_2(u(f),u(g))&=u(fg),&
 M_2(u(f),v(g))&=v(fg),&
 M_2(v(f),u(g))&=-v(fg).
\end{align*}
For \(k\ge2\), the remaining values are
\begin{equation}\label{eq:sc050-cups}
 M_k(v(g_1),\ldots,v(g_k))
       =u\left(\frac{P^{(k)}}{k!}\prod_i g_i\right).
\end{equation}
For \(k\ge3\), an input of type \(u\) gives zero.
Indeed the only possible outer operation is the scalar part
of \(b\), and its slots are counted exactly as in the proof
of \eqref{eq:sc050-braces}. These cases exhaust all cup operations.
Their higher identities are identities of actual insertions:
in an iterated insertion, partition terms by the ordered
positions of the inner blocks. Disjoint blocks occur in both
orders with their Koszul sign, and nested blocks give the
corresponding composite insertion. With \(b^2=0\), these
cancellations give the differential and higher cup identities.

In ordinary parity, the higher cup algebra has underlying module
\[
 \mathbb C[[t]]1\oplus\mathbb C[[t]]\eta,\qquad |\eta|=1,
\]
and
\begin{equation}\label{eq:sc050-Taylor}
 d\eta=P'(t),\qquad \eta\cdot\eta=-\tfrac12P''(t).
\end{equation}
Its higher suspended scalar operations are \(P^{(k)}(t)/k!\).
This is the Taylor family with its full normalized insertion
structure. The cup operations are \(\mathbb C[[t]]\)-linear.
The braces are only \(\mathbb C\)-linear: for example
\[
 v(t)\{v(1)\}=v(1),\qquad
 t\bigl(v(1)\{v(1)\}\bigr)=0.
\]
Thus adding insertion operations changes the required
coefficient-linearity of a comparison.

\subsection{Cohomology and detection by boundary words}

If \(P'\ne0\), multiplication by \(P'\) is injective on
\(\mathbb C[[t]]\). Equation \eqref{eq:sc050-HHd} therefore gives
\begin{equation}\label{eq:sc050-HH}
 HH^{\mathrm{even}}(H_P)=\mathbb C[[t]]/(P'),\qquad
 HH^{\mathrm{odd}}(H_P)=0.
\end{equation}
The notation \(HH\) denotes the cohomology of the full normalized
Hochschild complex just defined. The even product is the
quotient product, by the formula for \(M_2(u(f),u(g))\).
The induced odd bracket on cohomology is zero by parity.
The cochain braces \eqref{eq:sc050-braces} remain nonzero.

\subsection{The diagonal matrix for the full Taylor family}

The Taylor coefficients in \eqref{eq:sc050-cups} also have an
explicit matrix realization before setting their parameter to zero.
Rename the cochain parameter \(y\), put \(R=\mathbb C[[y]]\), and
let \(S=\mathbb C[[x,y]]\). These are rings of formal power series,
with convergence meaning agreement of every coefficient below a
fixed total degree. The series
\[
 g(x,y)=\frac{P(x)-P(y)}{x-y}
\]
belongs to \(S\): divide each monomial difference by its finite
geometric sum and add the resulting series. Define
\[
 \Delta_P=\begin{pmatrix}0&x-y\\g(x,y)&0\end{pmatrix},
 \qquad E_P=\operatorname{End}_S(S^{1|1}),\qquad
 \delta=[\Delta_P,-]_{\mathrm{gr}}.
\]
The square of \(\Delta_P\) is the central scalar \(P(x)-P(y)\),
so \(\delta^2=0\).

\begin{proposition}\label{prop:sc050-HH-diagonal}
The higher cup algebra of normalized Hochschild cochains of
\(H_P\), regarded as the finite free \(R\)-module in
\eqref{eq:sc050-Taylor}, has a strictly unital \(R\)-linear
\(A_\infty\)-quasi-isomorphism to \(E_P\).
Its components on \(e=s1\), \(a=s\eta\) are
\[
 \Psi_1(e)=sI,\qquad
 \Psi_n(a^{\otimes n})
      =-\frac{s\,\partial_y^n\Delta_P}{n!},\quad n\ge1,
\]
with \(x\) fixed under differentiation. The other
unit-containing components vanish.
\end{proposition}
\begin{proof}
Put \(f(t)=-(\Delta_P(x,y+t)-\Delta_P(x,y))\).
Subtracting the two squared matrices gives
\[
 \delta f-f^2=(P(y+t)-P(y))I.
\]
Its coefficient of \(t^n\) is the morphism identity for the
source operation \(P^{(n)}(y)e/n!\), including \(n=1\).
The unit equations have the same endpoint cancellations as
in Theorem~\ref{thm:sc050-boundary-matrix}. The cut-coproduct
induction proves the full morphism identity.

For its linear component set \(z_0=x-y\),
\(\beta=E_{21}\), and
\[
 H_0=-\partial_y\Delta_P=E_{12}-(\partial_yg)E_{21}.
\]
Then
\[
 \delta\beta=z_0I,\qquad \delta H_0=P'(y)I,\qquad
 H_0\beta+\beta H_0=I.
\]
The four matrices \(I,H_0,\beta,H_0\beta\) form an \(S\)-basis:
the odd pair is triangular in \(E_{12},E_{21}\), and the
even pair is \(I,E_{11}\).
Write \(S=R[[z_0]]\). Projection sets \(z_0=0\), kills terms
containing \(\beta\), and sends \(H_0\) to \(\eta\).
Inclusion sends \(1,\eta\) to \(I,H_0\).
For \(q\in R[[z_0]]\), let
\[
 \partial_0q=\frac{q(z_0)-q(0)}{z_0}.
\]
Define an odd homotopy by
\[
 h(qI)=(\partial_0q)\beta,\qquad
 h(qH_0)=-(\partial_0q)H_0\beta,\qquad
 h(q\beta)=h(qH_0\beta)=0.
\]
For the first basis element, \(\delta h(qI)=q-q(0)\).
For the second, the two terms proportional to
\(P'(y)(\partial_0q)\beta\) cancel, leaving
\((q-q(0))H_0\).
For \(q\beta\), the term \(h\delta(q\beta)=h(z_0qI)\)
is \(q\beta\).
For \(qH_0\beta\), the differential is
\(qP'(y)\beta-z_0qH_0\), and applying \(h\) gives
\(qH_0\beta\).
Thus \(\delta h+h\delta=1-ip\), and \(pi=1\).
The source differential is \(d\eta=P'(y)\), so these are
chain maps and the inclusion is a homotopy equivalence.
It is exactly \(\Psi_1\).

Each derivative and divided-difference operator lowers formal
order by a bounded amount. All specified components are
continuous. Their formulas use finite sums in each input
length and well-defined formal series in \(x,y\).
\end{proof}

This map compares the complete higher cup algebra with the
diagonal endomorphisms. It is coefficient-linear.
The insertions in \eqref{eq:sc050-braces} also differentiate
coefficients. A comparison that preserves those insertions
must include that additional, \(\mathbb C\)-linear structure;
it is not asserted by the matrix components \(\Psi_n\) alone.

\subsection{Recovering the scalar algebra from word evaluations}

For \(P=t^{r+1}/(r+1)\), every cohomology class has a unique
representative \(\sum_{j=0}^{r-1}f_jt^j\).
The cochain \(u(f)\) evaluated on \(j\) copies of the boundary
letter \(a\) gives \(f_je\).
These evaluations are unchanged by a boundary \(u(t^rg)\) for
\(0\le j<r\). Hence they give an isomorphism
\[
 HH^{\mathrm{even}}(H_P)\longrightarrow\mathbb C^r,\qquad
 [u(f)]\longmapsto(f_0,\ldots,f_{r-1}).
\]
They also recover the algebra product:
\[
 (fg)_n=\sum_{i+j=n}f_ig_j,\qquad 0\le n<r.
\]
In the cubic, the class \(t\) has zero zero-input value but
nonzero one-input value. In general \(t^{r-1}\) is first
detected by the \((r-1)\)-input word. The surviving scalar
boundary operation has arity \(r+1\); the two arities are
related by the explicit differential \(Dv(g)=u(t^rg)\).

The residue pairing is equally explicit. Define the residue of
a formal Laurent series times \(dt\) to be the coefficient of
\(t^{-1}dt\). Then
\[
 \lambda(f)=\operatorname{Res}_{t=0}\frac{f(t)\,dt}{t^r}
            =f_{r-1},\qquad
 \lambda(t^it^j)=\begin{cases}1&i+j=r-1,\\0&i+j\ne r-1.\end{cases}
\]
Its matrix in \(1,t,\ldots,t^{r-1}\) is antidiagonal with
nonzero entries one, so the pairing is nondegenerate.
This is a pairing on the calculated associative centre.
Identifying it with a chosen geometric or physical trace
requires a map that preserves the stated normalization.

\section{The polynomial retained by the differential}
\label{sec:sc050-charge}

The nilpotent differential of the fields has a short expression.
Put \(v(x)=x^r/(r+1)\) and
\[
 J=:x\bar\gamma:+:v(x)\gamma:.
\]
It is odd and has Hamiltonian weight one. Its zero mode sends
\(\gamma\) to \(x\), \(\bar\gamma\) to \(v(x)\), and \(x\) to
zero, so it is \(Q\). On the twisted module, multiplication of
the mode series shows that its zero mode is exactly
\(\mathcal Q\) in \eqref{eq:sc050-module-charge}.
Its self-bracket is
\[
 [J_\lambda J]=2xv(x)=2W(x).
\]
For an odd zero mode, Jacobi gives
\(Q^2=(J_{(0)}J)_{(0)}/2=W(x)_{(0)}\).
In \(V_r\), the field \(W(x)\) has only nonnegative powers of
position and acts by multiplication. Its zero mode vanishes.
This explains directly why \(Q^2=0\) in the constructed theory.
On \(\mathcal M\), its coefficient of \(z^{-1}\) is \(W(X)\).
The same field identity therefore gives the nonzero scalar
square in \eqref{eq:sc050-module-curvature}.

The same identity specifies what happens if the fields are
enlarged. Adjoin an even conjugate field \(p\) with
\([p_\lambda x]=1\), and keep \(J\) as written.
Skew symmetry gives \([x_\lambda p]=-1\), so the induced
differential must satisfy
\[
 Qp=-\bar\gamma-v'(x)\gamma,\qquad
 Q^2p=-v(x)-xv'(x)=-W'(x)=-x^r.
\]
It is not square zero on the enlarged algebra. A conjugate
momentum therefore requires a new differential or new
interacting fields with a proved cancellation of this term.
Adjoining a field \(c\) with \(Qc=1\) would instead make every
closed state exact, since \(Q(cA)=A\) when \(QA=0\).
That would eliminate the nontrivial cohomology calculated above.

There is no state whose zero mode implements translation on
\(V_r\). Such a state \(\omega\) would have to
satisfy \(\omega_{(0)}x=Tx\). But \(x\) has regular products
with every state, whereas \(Tx=x_1\ne0\). The Hamiltonian
grading and its zero-mode quotient remain well-defined.
They supply no additional gravitational fields or stress tensor.

The constructed maps preserve successively the singular vertex
operations, their higher conformal compatibility, and the
associative operations after the specified energy reduction.
The normalized Hochschild calculation retains every ordered
boundary insertion. A comparison with the full native chiral
centre must also retain nonzero spatial modes and their collision
operations. Neither the zero-mode quotient nor equality of the
scalar cohomology rings supplies that further map.
